\begin{lemma}[Proved: a cyclic Tate invariant for flat square-zero directions]
\label{C:cyclic_linear_Bockstein}
Let $A$ be a characteristic-zero DVR with maximal ideal $(p)$,
$F=A/p$, let $J$ be a $p$-torsion-free $A$-module, and put
$k=A\ltimes J$. Fix an $F$-linear functional $\ell:J/pJ\to F$.
For a finite free $k$-representation $M$ of $C_p$, there are natural
odd maps $b$ on the two Tate groups of $M_0=M\otimes_k A$.
For every monic irreducible $q\in F[z]$ other than $z$,
\[
 \chi_{q,\ell}(M)=
 \operatorname{length}_{F[z]}
   \bigl(\widehat H^0(C_p,M_0),z=b^2\bigr)_q\pmod2
\]
is additive on admissible exact sequences.
\end{lemma}
\begin{proof}
The $A$-module $J$ is flat. Choose a basis of $M$ and write its
generator as $T=T_0+U$, with entries of $U$ in $J$.
The two-periodic Tate differential is
$d^0=T-I$, $d^1=I+T+\cdots+T^{p-1}$.
Write $d=d_0+b_J$ according to $k=A\oplus J$.
The identity $d^2=0$ gives
$d_0b_J+b_Jd_0=0$. Hence $b_J$ induces an odd map from
$H(d_0)$ to $H(d_0\otimes_AJ)$.
Flatness gives the natural identification
\[
 H(d_0\otimes_AJ)=H(d_0)\otimes_AJ
              =H(d_0)\otimes_F(J/pJ),
\]
since both Tate groups are annihilated by $p$.
Compose with $\ell$ to obtain $b$.
For a map $f=f_0+f_J$ between two such complexes, the coefficient
of $J$ in $d_Nf=fd_M$ says
$b_{J,N}f_0-f_0b_{J,M}=f_Jd_{0,M}-d_{0,N}f_J$.
It follows that $b$ is natural and independent of the chosen basis.

For an admissible sequence, choose an underlying $k$-linear splitting.
Write its differentials as
$d_E^i=\left(\begin{smallmatrix}d_A^i&X^i\\0&d_B^i\end{smallmatrix}\right)$,
and write $X^i=X_0^i+X_J^i$.
The $J$-coefficient of $d_E^{i+1}d_E^i=0$ is
\[
 d_{0,A}^{i+1}X_J^i+b_{J,A}^{i+1}X_0^i
       +X_0^{i+1}b_{J,B}^i+X_J^{i+1}d_{0,B}^i=0.
\]
On homology this gives $b\partial+\partial b=0$ for the connecting
map of the exact Tate six-cycle. Thus every map of that cycle
commutes with $z=b^2$, and primary decomposition remains exact.
On a primary part with $q\ne z$, the odd map $b$ is invertible.
Use it to identify the odd Tate group of each term with its even
Tate group. The six-cycle becomes an exact cyclic sequence on the
three even groups: the second connecting map differs after folding
by the invertible operator $-z$, which preserves the relevant
kernels and images. If $r,s,t$ are the lengths of its three images,
the lengths of its terms are $r+t,r+s,s+t$. Their alternating
difference is even, proving the invariant.
\end{proof}

\begin{proposition}[Proved: all unramified flat square-zero directions]
\label{C:unramified_flat_squarezero_allprimes}
Under the preceding assumptions,
\[
 \vartheta_{A\ltimes J}\text{ is an isomorphism in every degree}
       \quad\Longleftrightarrow\quad pJ=J.
\]
If $J/pJ\ne0$, failure occurs already in degree $p$, by an explicit
representation of rank $p!$ induced from a rank-$p$ representation
of $C_p$.
\end{proposition}
\begin{proof}
If $pJ=J$, multiplication by $p$ is bijective and all integers prime
to $p$ are already units in $A$. Thus $J$ is uniquely divisible by
every positive integer. Proposition~\ref{C:divisible_nil_invariance}
and Proposition~\ref{C:DVR} prove sufficiency.

Suppose $J/pJ\ne0$, choose $x\in J$ and $\ell$ with $\ell(\bar x)=1$.
Let
\[
 L=A[t]/(1+t+\cdots+t^{p-1}),\qquad C:L\to L,
       \quad C(y)=ty.
\]
In the basis $1,t,\ldots,t^{p-2}$, let $b:A\to L$ send $1$ to $1$,
and let $a:L\to A$ extract the coefficient of $t^{p-2}$.
On $k\oplus(L\otimes_Ak)$ put
\[
 T=T_0+xU,\qquad
 T_0=\begin{pmatrix}1&0\\0&C\end{pmatrix},\qquad
 U=\begin{pmatrix}0&a\\b&0\end{pmatrix}.
\]
Since $J^2=0$, the coefficient of $x$ in $T^p$ has off-diagonal
blocks $a\sum_{i=0}^{p-1}C^i$ and
$\sum_{i=0}^{p-1}C^ib$, both zero. Thus $T^p=1$.
The even and odd Tate groups of $A\oplus L$ are respectively
$A/p$ (from $A$) and $L/(C-1)L=A/p$ (from $L$).
The first odd map induced by $U$ sends the class of $1$ to that of
$b(1)=1$.
For the other map, the derivative of the norm has upper-right block
$aQ(C)$, where
\[
 Q(t)=\sum_{i=0}^{p-2}(p-1-i)t^i.
\]
Its value on $1$ is $aQ(C)b=1$ by the choice of the last-coefficient
functional. Consequently both odd maps are $1$, so $b^2=1$ and
$\chi_{z-1,\ell}(M)=1$.

To retain this detector after induction to $\mathfrak S_p$, twist
the generator by $g\mapsto g^r$, $1\le r<p$.
The norm is unchanged, since multiplication by $r$ permutes the
exponents modulo $p$ and $T^p=1$. The derivative of $T^r-I$ sends
$1$ to $\sum_{i=0}^{r-1}C^ib$, whose class in
$L/(C-1)L$ is $r$. The ideals $(C^r-1)L$ and $(C-1)L$ agree:
the two generators of $C_p$ generate the same cyclic group, and
the usual geometric-sum identities give both inclusions.
Thus $b^2=r$ on both one-dimensional Tate groups of the twist.
Exactly the twist $r=1$ has a nonzero $(z-1)$-primary part.

In the Mackey restriction of
$\operatorname{Ind}_{C_p}^{\mathfrak S_p}M$, the double cosets in
the normalizer give these $p-1$ twists once each. Every other
intersection of two distinct order-$p$ subgroups is trivial;
its summand is induced from the trivial subgroup and has zero
Tate groups. Hence
\[
 \chi_{z-1,\ell}\!\left(
  \operatorname{Res}_{C_p}^{\mathfrak S_p}
  \operatorname{Ind}_{C_p}^{\mathfrak S_p}M\right)=1.
\]
Every representation obtained by base change from $A$ has zero
odd map, so every Young permutation representation has invariant
zero, also after restriction. This proves failure of degree-$p$
surjectivity. The induced representation has rank
$p[\mathfrak S_p:C_p]=p!$, as asserted.
\end{proof}

The finite certificate \texttt{code/C\_finish\_cyclic\_flat\_probe.py} checks the actual integral cyclotomic matrices, Tate kernels and norm images, Smith invariants, and all $23$ normalizer twists for $p=2,3,5,7,11$. Its output is \texttt{logs/C\_finish\_cyclic\_flat\_probe.json}; finite matrix checks supplement the arbitrary-module and exact-sequence proof above.
